\documentclass[10pt]{article}

% ============================================================
% Packages
% ============================================================
\usepackage[letterpaper,margin=0.85in]{geometry}
\usepackage{amsmath,amssymb,amsfonts}
\usepackage{amsthm}
\usepackage{mathtools}
\usepackage{microtype}
\usepackage{fancyhdr}
\usepackage[hidelinks]{hyperref}

% ============================================================
% Typography and spacing
% ============================================================
\setlength{\parindent}{1.2em}
\setlength{\parskip}{0pt}
\setlength{\textfloatsep}{10pt}
\setlength{\abovedisplayskip}{7pt}
\setlength{\belowdisplayskip}{7pt}

% ============================================================
% Theorem environments
%
% Separate counters intentionally match the conference paper:
% Lemma 1, Theorem 1, Corollary 1, Corollary 2.
% ============================================================
\newtheorem{theorem}{Theorem}
\newtheorem{lemma}{Lemma}
\newtheorem{corollary}{Corollary}

% ============================================================
% Header / footer
% ============================================================
\pagestyle{fancy}
\fancyhf{}
\fancyhead[L]{\small Mishal and Koyuncu}
\fancyhead[R]{\small LR-PPO Technical Report}
\fancyfoot[C]{\thepage}
\renewcommand{\headrulewidth}{0.4pt}
\renewcommand{\footrulewidth}{0pt}

% ============================================================
% Document
% ============================================================
\begin{document}

\begin{center}

{\LARGE\bfseries
Hybrid Predictive Guidance with Residual Reinforcement Learning\\
for Autonomous UAS Interception
\par}

\vspace{0.6em}

{\large\bfseries
Technical Report: Proofs of Theoretical Results
\par}

\vspace{1em}

{\large
Michael Mishal \qquad Erdem Koyuncu
\par}

\vspace{0.25em}

{\normalsize
Department of Electrical and Computer Engineering\\
University of Illinois Chicago, Chicago, IL, USA\\
\texttt{\{mmisha2, ekoyuncu\}@uic.edu}
\par}

\vspace{0.6em}

{\small August 2026}

\end{center}

\vspace{0.6em}

\begin{abstract}
This technical report supplements the paper
\emph{Hybrid Predictive Guidance with Residual Reinforcement Learning for
Autonomous UAS Interception}. It provides detailed proofs of the geometric
properties of the proposed Lead-Residual Proximal Policy Optimization
(LR-PPO) guidance law. Specifically, the results establish orthogonality and
norm bounds for the projected residual, an explicit bound on angular
deviation from predictive Lead guidance, exact recovery of the nominal Lead
direction for residuals with no orthogonal component, and a positive lower
bound on alignment between the executed LR-PPO and Lead commands. These
properties constrain the authority of the learned correction but do not imply
closed-loop stability or global interception optimality.
\end{abstract}


% ============================================================
% Theoretical statements
% ============================================================
\section*{Theoretical Properties}
\label{sec:theoretical-results}

Let \(u_t^{\mathrm{Lead}}\) denote the unit nominal Lead guidance direction.
The PPO policy produces a raw residual \(r_t^{\mathrm{raw}}\), which is clipped
to the unit ball:
\[
r_t=
\begin{cases}
\dfrac{r_t^{\mathrm{raw}}}{\|r_t^{\mathrm{raw}}\|},
& \|r_t^{\mathrm{raw}}\|>1,\\[4pt]
r_t^{\mathrm{raw}},
& \text{otherwise}.
\end{cases}
\]
Thus,
\[
\|r_t\|\leq1.
\]

The component of this residual orthogonal to Lead is
\[
r_t^\perp
=
r_t-
\left(r_t^T u_t^{\mathrm{Lead}}\right)
u_t^{\mathrm{Lead}},
\]
and the executed LR-PPO direction is
\[
u_t^{\mathrm{LR}}
=
\operatorname{normalize}
\left(
u_t^{\mathrm{Lead}}
+
\tan(\theta_{\max})r_t^\perp
\right),
\]
where
\[
\operatorname{normalize}(x)
=
\frac{x}{\|x\|},
\qquad
0<\theta_{\max}<\frac{\pi}{2}.
\]
The controller evaluated in the accompanying paper uses
\(\theta_{\max}=30^\circ\).

\begin{lemma}
\label{lem:orthogonal_residual}
Let \(u_t^{\mathrm{Lead}}\) be a unit Lead direction and let the clipped
residual satisfy \(\|r_t\|\leq1\). Then
\[
r_t^\perp
\perp
u_t^{\mathrm{Lead}},
\qquad
\|r_t^\perp\|\leq1.
\]
\end{lemma}

\begin{theorem}
\label{thm:angular_bound}
For
\[
u_t^{\mathrm{LR}}
=
\operatorname{normalize}
\left(
u_t^{\mathrm{Lead}}
+
\tan(\theta_{\max})r_t^\perp
\right),
\]
with \(0<\theta_{\max}<\pi/2\),
\[
\angle
\left(
u_t^{\mathrm{LR}},
u_t^{\mathrm{Lead}}
\right)
\leq
\theta_{\max}.
\]
\end{theorem}

\begin{corollary}
\label{cor:lead_recovery}
If \(r_t^\perp=0\), then
\[
u_t^{\mathrm{LR}}
=
u_t^{\mathrm{Lead}}.
\]
Thus, zero, parallel, or anti-parallel residuals recover the nominal Lead
direction exactly.
\end{corollary}

\begin{corollary}
\label{cor:positive_alignment}
The LR-PPO command always retains positive alignment with Lead:
\[
\left(u_t^{\mathrm{LR}}\right)^T
u_t^{\mathrm{Lead}}
\geq
\cos(\theta_{\max})
>
0.
\]
For \(\theta_{\max}=30^\circ\),
\[
\cos(30^\circ)
=
\frac{\sqrt{3}}{2}
\approx0.866.
\]
\end{corollary}

These results constrain the learned controller geometrically: PPO may modify
the predictive Lead command within the prescribed angular neighborhood, but
it cannot arbitrarily overwrite or reverse the nominal guidance direction.


% ============================================================
% Proofs
% ============================================================
\section*{Proofs}

\subsection*{Proof of Lemma~\ref{lem:orthogonal_residual}}

Recall that
\begin{equation}
r_t^\perp
=
r_t-
\left(r_t^T u_t^{\mathrm{Lead}}\right)
u_t^{\mathrm{Lead}}.
\label{eq:app_residual_projection}
\end{equation}
Since
\(
\|u_t^{\mathrm{Lead}}\|=1
\),
taking the inner product with the Lead direction gives
\begin{align}
\left(u_t^{\mathrm{Lead}}\right)^T r_t^\perp
&=
\left(u_t^{\mathrm{Lead}}\right)^T r_t
-
\left(r_t^T u_t^{\mathrm{Lead}}\right)
\left(u_t^{\mathrm{Lead}}\right)^T
u_t^{\mathrm{Lead}} \\
&=
r_t^T u_t^{\mathrm{Lead}}
-
r_t^T u_t^{\mathrm{Lead}} \\
&=0.
\end{align}
Therefore,
\begin{equation}
r_t^\perp
\perp
u_t^{\mathrm{Lead}}.
\label{eq:app_orthogonality}
\end{equation}

The residual can be decomposed into mutually orthogonal components:
\begin{equation}
r_t
=
r_t^\perp
+
\left(r_t^T u_t^{\mathrm{Lead}}\right)
u_t^{\mathrm{Lead}}.
\end{equation}
By the Pythagorean relation,
\begin{equation}
\|r_t\|^2
=
\|r_t^\perp\|^2
+
\left(r_t^T u_t^{\mathrm{Lead}}\right)^2.
\end{equation}
Hence,
\begin{equation}
\|r_t^\perp\|^2
\leq
\|r_t\|^2.
\end{equation}
Since the residual is clipped to the unit ball,
\(
\|r_t\|\leq1
\),
so
\begin{equation}
\|r_t^\perp\|
\leq
\|r_t\|
\leq1.
\end{equation}
This establishes both the orthogonality and norm-bound claims.


\subsection*{Proof of Theorem~\ref{thm:angular_bound}}

Define the unnormalized LR-PPO command
\begin{equation}
\tilde u_t^{\mathrm{LR}}
=
u_t^{\mathrm{Lead}}
+
\tan(\theta_{\max})r_t^\perp.
\label{eq:app_lr_command}
\end{equation}
The executed command is
\begin{equation}
u_t^{\mathrm{LR}}
=
\frac{\tilde u_t^{\mathrm{LR}}}
{\|\tilde u_t^{\mathrm{LR}}\|}.
\label{eq:app_lr_normalized}
\end{equation}

Let
\begin{equation}
\phi_t
=
\angle
\left(
u_t^{\mathrm{LR}},
u_t^{\mathrm{Lead}}
\right).
\end{equation}
Normalization changes only magnitude, so \(\phi_t\) is also the angle between
\(\tilde u_t^{\mathrm{LR}}\) and \(u_t^{\mathrm{Lead}}\).

Using Lemma~\ref{lem:orthogonal_residual},
\begin{align}
\left(u_t^{\mathrm{Lead}}\right)^T
\tilde u_t^{\mathrm{LR}}
&=
\left(u_t^{\mathrm{Lead}}\right)^T
u_t^{\mathrm{Lead}}
+
\tan(\theta_{\max})
\left(u_t^{\mathrm{Lead}}\right)^T
r_t^\perp \\
&=1.
\end{align}
Thus,
\(\tilde u_t^{\mathrm{LR}}\neq0\) and
\[
0\leq\phi_t<\frac{\pi}{2}.
\]

Equation~\eqref{eq:app_lr_command} consists of a unit component along
\(u_t^{\mathrm{Lead}}\) and an orthogonal component of magnitude
\[
\tan(\theta_{\max})\|r_t^\perp\|.
\]
Therefore,
\begin{equation}
\tan(\phi_t)
=
\tan(\theta_{\max})\|r_t^\perp\|.
\label{eq:app_angle_relation}
\end{equation}
Lemma~\ref{lem:orthogonal_residual} gives
\(
\|r_t^\perp\|\leq1
\),
hence
\[
\tan(\phi_t)
\leq
\tan(\theta_{\max}).
\]
Because tangent is strictly increasing on
\([0,\pi/2)\),
\[
\phi_t
\leq
\theta_{\max}.
\]
Consequently,
\begin{equation}
\angle
\left(
u_t^{\mathrm{LR}},
u_t^{\mathrm{Lead}}
\right)
\leq
\theta_{\max}.
\end{equation}

For the controller used in the accompanying paper,
\(\theta_{\max}=30^\circ\); therefore the executed LR-PPO direction can deviate
by at most \(30^\circ\) from the nominal Lead direction.


\subsection*{Proof of Corollary~\ref{cor:lead_recovery}}

Suppose
\[
r_t^\perp=0.
\]
Then from \eqref{eq:app_lr_command},
\[
\tilde u_t^{\mathrm{LR}}
=
u_t^{\mathrm{Lead}}.
\]
Because the Lead direction is unit norm,
\[
u_t^{\mathrm{LR}}
=
\frac{u_t^{\mathrm{Lead}}}
{\|u_t^{\mathrm{Lead}}\|}
=
u_t^{\mathrm{Lead}}.
\]

In particular, if the clipped residual is parallel to Lead,
\[
r_t
=
\alpha u_t^{\mathrm{Lead}},
\qquad
|\alpha|\leq1,
\]
then
\begin{align}
r_t^\perp
&=
\alpha u_t^{\mathrm{Lead}}
-
\left(
\alpha
\left(u_t^{\mathrm{Lead}}\right)^T
u_t^{\mathrm{Lead}}
\right)
u_t^{\mathrm{Lead}} \\
&=
\alpha u_t^{\mathrm{Lead}}
-
\alpha u_t^{\mathrm{Lead}} \\
&=0.
\end{align}
Clipping preserves parallel or anti-parallel orientation. Therefore, a raw PPO
residual that is zero, parallel, or anti-parallel to the Lead direction also
produces \(r_t^\perp=0\), recovering the nominal Lead command exactly.


\subsection*{Proof of Corollary~\ref{cor:positive_alignment}}

Because
\(u_t^{\mathrm{LR}}\)
and
\(u_t^{\mathrm{Lead}}\)
are unit vectors,
\begin{equation}
\left(u_t^{\mathrm{LR}}\right)^T
u_t^{\mathrm{Lead}}
=
\cos(\phi_t),
\end{equation}
where
\[
\phi_t
=
\angle
\left(
u_t^{\mathrm{LR}},
u_t^{\mathrm{Lead}}
\right).
\]

By Theorem~\ref{thm:angular_bound},
\[
0\leq\phi_t\leq\theta_{\max}<\frac{\pi}{2}.
\]
Since cosine is decreasing on this interval,
\[
\cos(\phi_t)
\geq
\cos(\theta_{\max}).
\]
Therefore,
\begin{equation}
\left(u_t^{\mathrm{LR}}\right)^T
u_t^{\mathrm{Lead}}
\geq
\cos(\theta_{\max})
>
0.
\end{equation}

For \(\theta_{\max}=30^\circ\),
\begin{equation}
\left(u_t^{\mathrm{LR}}\right)^T
u_t^{\mathrm{Lead}}
\geq
\cos(30^\circ)
=
\frac{\sqrt{3}}{2}
\approx0.866.
\end{equation}
Hence, LR-PPO always retains a positive directional component along the
nominal Lead command and cannot reverse the analytical guidance direction.


% ============================================================
% Closing note
% ============================================================
\section*{Scope of the Results}

The preceding results describe geometric properties of the LR-PPO action
construction. They guarantee bounded deviation from predictive Lead guidance,
exact nominal-controller recovery for residuals without an orthogonal
component, and positive alignment with the Lead direction. They do not
constitute a proof of closed-loop stability, global interception optimality,
or guaranteed interception under arbitrary target motion.


\end{document}